% Generated draft: requires mathematical and source review.
\documentclass[11pt]{article}
\usepackage[margin=25mm]{geometry}
\usepackage{fontspec}
\setmainfont{Latin Modern Roman}
\usepackage{amsmath,amssymb,mathtools,unicode-math}
\setmathfont{Latin Modern Math}
\usepackage{enumitem,needspace,xurl,hyperref}
\setlength{\parindent}{0pt}
\setlength{\parskip}{5pt}
\setlength{\emergencystretch}{3em}
\newcommand{\N}{\mathbb{N}}
\newcommand{\Z}{\mathbb{Z}}
\newcommand{\Q}{\mathbb{Q}}
\newcommand{\R}{\mathbb{R}}
\newcommand{\C}{\mathbb{C}}
\newcommand{\F}{\mathbb{F}}
\newcommand{\A}{\mathbb{A}}
\newcommand{\PP}{\mathbb{P}}
\newcommand{\E}{\mathbb{E}}
\newcommand{\eps}{\varepsilon}
\newcommand{\dd}{\,\mathrm{d}}
\begin{document}
\section*{UnsolvedMath ID: 7200008}
\textbf{Source problem code:} catalogue entry
\subsection{Definitions and mathematical statement}
In the standard threefold terminology, an elephant is an effective anticanonical divisor \(D\in|-K_X|\). For a normal complex threefold with \(K_X\) \(\Q\)-Cartier, terminal singularities mean that every exceptional divisorial discrepancy \(a(E,X)\) is positive. A normal complex surface has at most Du Val singularities if every singular point is a rational double point of type \(A\), \(D\), or \(E\); smooth points are allowed.

The supplied question is: do general elephants have at most Du Val singularities? Once an intended class of threefolds \(X\) and a nonempty anticanonical system are specified, the usual global formulation asks whether there is a dense Zariski open subset
\[
 U\subseteq|-K_X|=\mathbb P H^0(X,\mathcal O_X(-K_X))
\]
whose members are normal surfaces with at most Du Val singularities. Here \(\mathcal O_X(-K_X)\) is understood reflexively when \(-K_X\) is only a Weil divisor.

The excerpt does not specify the class of \(X\), base field, or whether the question is local near a terminal singularity, global for projective threefolds, or restricted to \(\Q\)-Fano threefolds. Nonemptiness of \(|-K_X|\) is also absent. The supplied summary speaks broadly of terminal threefolds and cites work on \(\Q\)-Fano threefolds and terminal flips; those distinct settings must not be collapsed into an unconditional assertion for every terminal threefold. The displayed formulation explains the usual global interpretation, with the missing hypotheses still unresolved.
\subsection{Short English statement}
Does a general anticanonical surface section of the intended threefold have only Du Val singularities? The source does not identify the precise local or global threefold class or guarantee that the anticanonical system is nonempty.
\subsection{Sources}
\begin{itemize}
\item Local UnsolvedMath registry, permanent ID 7200008, source code catalogue entry, updated 2026-08-06; no proposer or proposal year supplied.
\item Wikipedia, \emph{List of unsolved problems in mathematics}, Geometry bullet 8, pinned revision 1366636767: \url{https://en.wikipedia.org/w/index.php?title=List_of_unsolved_problems_in_mathematics&oldid=1366636767#Geometry}. Registry access date 2026-07-29; MediaWiki extraction. Wikipedia text is CC BY-SA 4.0; attribution and ShareAlike handling remain required for release. The provenance says the extracted bullet remains unpublished pending primary-source status review, despite the registry's published flag.
\item Supplied research-summary references: arXiv:1404.0909, ``Deforming elephants of Q-Fano threefolds''; arXiv:1608.00364, ``Normality of general elephants on 3-fold terminal flips''. No author names or full bibliographic entries are supplied. The summary's reported searches and August 2026 open-status assessment are not independent checks performed for this draft.
\end{itemize}
\end{document}
